\documentclass[conference]{IEEEtran}
\IEEEoverridecommandlockouts
\usepackage{cite}
\usepackage{amsmath,amssymb,amsfonts}
\usepackage{algorithmic}
\usepackage{algorithm}
\usepackage{graphicx}
\usepackage{textcomp}
\usepackage{xcolor}
\usepackage{hyperref}
\usepackage{booktabs}
\usepackage{flushend} % Balances the columns on the final page

\def\BibTeX{{\rm B\kern-.05em{\sc i\kern-.025em b}\kern-.08em
    T\kern-.1667em\lower.7ex\hbox{E}\kern-.125emX}}


\begin{document}
\title{Blockchain Based Farm-to-Fork Food Traceability System With AI/ML Integrations}
\author{
\begin{tabular}{ccc}

Bhagyashree Ambore
&
Tanush Patil J
&
Vellanki Sricharan
\\
\textit{Dept. of CSE}
&
\textit{Dept. of ISE}
&
\textit{Dept. of ISE}
\\
RNS Institute of Technology
&
RNS Institute of Technology
&
RNS Institute of Technology
\\
Bengaluru, India
&
Bengaluru, India
&
Bengaluru, India
\\
ambore.bhagyashree@gmail.com
&
tanushpatilj23is@rnsit.ac.in
&
vellankisricharan23is@rnsit.ac.in
\\[2em]

B S Krishna
&
&
Supreet Kumar
\\
\textit{Dept. of ISE}
&
&
\textit{Dept. of ISE}
\\
RNS Institute of Technology
&
&
RNS Institute of Technology
\\
Bengaluru, India
&
&
Bengaluru, India
\\
bskrishna23is@rnsit.ac.in
&
&
supreetkumar24isdip@rnsit.ac.in

\end{tabular}
}
\maketitle
\begin{abstract}
At the moment, tracking agricultural produce to its origin is very large challenge as there is lack of market information transparency. This loss of information is the opportunity that middlemen take advantage of and reduces the power of farmers in their negotiations over their produce. To address this, we created FarmChain AI - a platform specifically created to empower farmers at harvest. We integrated Python-powered AI predictive tools with smart contracts on Ethereum that are untamperable. The heart of the system is built with React 18, Node.js, Express, Supabase, FastAPI and Solidity. The platform implements the complete supply chain lifecycle across six stakeholder roles: Farmer, Admin, Processor, Distributor, Retailer, and Consumer. Each custody transfer is permanently recorded on the Ethereum Sepolia Testnet via a \texttt{FarmChain.sol} smart contract that enforces cryptographic ownership verification. For market price prediction, we trained an XGBoost Regressor on 2{,}940 historical price records spanning 7 commodities across 7 Indian agricultural markets. The model uses 30 leakage-free time-series features and achieves a MAPE of 2.72\% and $R^2$ of 0.984 on a chronological test split of 294 samples. Data retrieval is free, but data posting requires a small gas fee ($\approx 0.00037$ ETH) to be placed on the blockchain. Consumer verification is facilitated through QR-code scanning and direct on-chain reads, providing a six-milestone authenticity checklist. At 14ms of API response time and 35ms for the price prediction of AI, this configuration shows that an affordable, responsive, and transparent harvest-logging software system is achievable without putting farmers under high pressure to pay expensive prices for enterprise corporate software.
\end{abstract}

\begin{IEEEkeywords}
Blockchain, Ethereum Sepolia, Smart Contracts, Agricultural Technology, FastAPI, Solidity, React.
\end{IEEEkeywords}

\section{Introduction}

\subsection{Background and Motivation}
In developing agrarian countries like the India, agriculture is the backbone of the country's economy with 50.3\% of people employed in agriculture and contributing to 17.5\% of total GDP of the country [1]. Regardless of this vast economic impact, transport from farm-to-market and from timber landing to the initial logging sites are still very disorganized. Farmers have limited knowledge of market demands, and frequently they have no objective evidence of their crops' worth, based on data. This is a structural asymmetry of information, and often they leave the negotiation table with only a quarter to a third of the ultimate selling price [2]. 

Traditional tracking software is based on single-company, central databases. When one corporate entity is controlling the server, it is easy to manipulate harvest records from the past and thereby break the trust and integrity of the data. Massive enterprise blockchains (like IBM Food Trust) do solve this issue of centralization, but at an unaffordable price tag for small independent farm cooperatives [9].

FarmChain AI is developed to give the farmer data sovereignty and negotiating power back. Our system leverages targeted AI predictive endpoints and public smart contracts, none of which charge any subscription fees; therefore, our system is decentralized and avoids the expense of having to pay for any corporate software.

\subsection{System Architecture Scope \& Tech Stack}
We utilized a decoupled open-source tech stack to architect the platform:
\begin{itemize}
    \item \textbf{Frontend Layer:} The application is a React 18 application styled with Tailwind CSS. The platform implements role-specific dashboards for all six stakeholder roles: Farmer, Admin, Processor, Distributor, Retailer, and Consumer.
    \item \textbf{Backend API \& Database:} A node.js and express based backend system routes and authenticates the transaction; connected to Supabase with a PostgreSQL database to store non-critical information off-chain.
    \item \textbf{AI Microservice:} The key AI functionality is encapsulated in a Python script. An XGBoost-based price prediction pipeline is hosted using FastAPI \cite{dragoni2017microservices}.
    \item \textbf{Blockchain:} The blockchain is a smart contract called \texttt{FarmChain.sol} written in native Solidity 0.8.24 and deployed on public Sepolia testnet of Ethereum \cite{wood2014ethereum}.
    \item \textbf{Event Streaming:} Apache Kafka with a transactional outbox pattern handles asynchronous event propagation across services.
\end{itemize}

\subsection{Paper Organization}
In the rest of this document, we give a comprehensive description of our approach. The previous research and literature in the domain is covered in Section II. In Section III, we will be introducing our four-layer technical architecture and the multi-stakeholder workflow. Section IV presents the AI/ML price prediction pipeline including dataset, features, model training, and evaluation. A detailed explanation of the coding implementation is provided in Section V. Our empirical speed and gas cost testing is presented in Section VI. Last but not least, in Sections VII and VIII, we recap our main results and give a vision of our future goals.

\section{Related Work}

\subsection{Blockchain in Agricultural Systems}
One of the first works to propose the concept of incorporating tracking tags with distributed ledgers for food items was done by Tian et al. \cite{tian2017agrifood}. Lin et al. \cite{lin2019bifts} further introduced the BIFTS system over Hyperledger. Unfortunately, the expensive hardware needed of these enterprise systems was not applicable to average farmers. The amazing cryptographic security that blockchains offer is a thing of beauty, but the simplicity and usability of web front-ends to the technology are a thing of desperate need for farm workers, as noted by Kamath \cite{kamath2018food} in their examination of Walmart's corporate Ethereum pilots. Antonucci et al. \cite{antonucci2019review} also noted that usability is a significant challenge for blockchain use in rural areas.

\subsection{AI Microservices in Agriculture}
On the topic of AI in agriculture, Kamilaris et al. \cite{kamilaris2018deep} described in detail a range of special-purpose machine learning services that can be applied to devices in the agricultural edge. In addition, Liakos et al. \cite{liakos2018ml} and Xiong et al. \cite{xiong2020price} have demonstrated that it is extremely useful to study historical market data patterns to predict future prices for crops. Chen and Guestrin \cite{chen2016xgboost} introduced XGBoost, which has since become the dominant algorithm for tabular time-series regression tasks in agricultural forecasting. In our project, FarmChain AI, we bring these intensive AI computations for price forecasting directly to the farmer. How we do this is to split the bulk of the math away from the easy-to-interact farmer dashboard into a small, easily accessed, but powerful FastAPI service that communicates asynchronously with the farmer service through REST calls.

\subsection{Integrated Systems \& Research Gap}
Although some projects are trying to leverage public Ethereum ledgers, a full-range platform that integrates AI market predictions with rigorous evaluation, multi-stakeholder custody tracking, consumer verification, and a genuinely farmer-friendly UI is missing. Existing solutions are geared towards enterprise distributors. This project is a complete, working prototype custom-designed to meet the technological gap at the farm gate.

\section{System Architecture and Design}

\subsection{Layered System Architecture}
We've broken our ecosystem into four highly specialized layers in our structural design:

\begin{figure}[htbp]
\centering
\includegraphics[width=0.95\linewidth]{architecture.png}
\caption{The System Architecture of FarmChain AI showcases the Presentation Layer, Express/Supabase API, FastAPI AI Microservices, and Blockchain Layer.}
\label{fig:architecture}
\end{figure}

\begin{enumerate}
    \item \textbf{Presentation Layer (React 18):} Role-specific dashboards for Farmer, Processor, Distributor, Retailer, Admin, and Consumer.
    \item \textbf{Backend API Layer (Node.js):} A light backend server to process user logins via JWT, role-gated routing, and manage database interactions.
    \item \textbf{AI Microservice Layer (FastAPI):} A safe Python environment containing our XGBoost price prediction pipeline.
    \item \textbf{Blockchain Persistence Layer (Solidity):} A permanently logged and tamper-proof record of all harvest data, stage updates, and ownership transfers on the Sepolia network.
\end{enumerate}

\subsection{Multi-Stakeholder Workflow Design}
The end-to-end custody flow proceeds through the following stages:
\begin{enumerate}
    \item \textbf{Farmer Registration:} The farmer logs harvest data (crop type, weight, expiry, photo). The system generates a batch ID (\texttt{BCH-YYYYMMDD-XXXX}) and fetches an AI price prediction from FastAPI. Status: \textit{Pending Approval}.
    \item \textbf{Admin Approval \& On-Chain Commitment:} An administrator reviews the submission and triggers \texttt{registerProduct()} on the smart contract via a backend relayer wallet. The blockchain transaction hash is stored in Supabase.
    \item \textbf{Processor Intake:} The processor browses available harvests and expresses interest. Upon farmer dispatch, the processor performs a \textbf{gate-scale weight verification} against the farmer's declared weight, followed by a \textbf{physical quality inspection} (Appearance, Freshness, Cleanliness, Grade A/B/C). Accepted batches proceed through processing stages; rejected batches are returned with documented reasons. The stage is updated on-chain via \texttt{updateStage()}.
    \item \textbf{Distributor Custody:} The distributor takes custody via \texttt{transferOwnership()}, which reassigns the \texttt{currentOwner} field on-chain. The distributor manages cold-chain transit and delivery logistics.
    \item \textbf{Retailer Receipt:} The retailer scans a QR code to verify the batch, then records receipt on-chain. The retailer module tracks inventory, point-of-sale transactions, wastage, expiry, and returns.
    \item \textbf{Consumer Verification:} Any consumer can enter a Product ID or scan a QR code on the public \texttt{/track} page. The system performs \textbf{direct on-chain reads} via \texttt{ethers.JsonRpcProvider} and presents a six-milestone authenticity checklist.
\end{enumerate}

\subsection{Smart Contract Architecture}
The \texttt{FarmChain.sol} contract (Solidity 0.8.24) stores two structures: a \texttt{Product} struct containing the ID, name, category, batch number, quantity, farmer address, current owner, status, and timestamp; and a \texttt{HistoryEvent} struct that logs each custody change. Four public functions constitute the contract interface:
\begin{itemize}
    \item \texttt{registerProduct()}: Creates a new on-chain record with the caller as farmer and initial owner. Emits \texttt{ProductRegistered}.
    \item \texttt{updateStage()}: Updates the product status. Enforces \texttt{require(currentOwner == msg.sender)} \cite{atzei2017survey}. Emits \texttt{StageUpdated}.
    \item \texttt{transferOwnership()}: Reassigns \texttt{currentOwner} to a new wallet address. Emits \texttt{OwnershipTransferred}.
    \item \texttt{getProductHistory()}: A gas-free \texttt{view} function returning the complete provenance audit trail.
\end{itemize}

\subsection{Off-Chain Input Validation}
Since blockchain immutability cannot guarantee that initially entered data are truthful, FarmChain AI implements three complementary validation mechanisms:
\begin{enumerate}
    \item \textbf{Admin Approval Gate:} Products enter as \textit{Pending Approval}. Only after administrative review does the system commit the record on-chain, preventing unverified data from reaching the immutable ledger.
    \item \textbf{Processor Physical Inspection:} The processor performs a gate-scale weight check (cross-referencing the farmer's declared weight) and a multi-criteria quality inspection (Appearance, Freshness, Cleanliness) with Accept/Reject decisions.
    \item \textbf{AI Anomaly Detection:} The AI microservice includes a fraud-scoring endpoint that flags statistically improbable input combinations before blockchain commitment.
\end{enumerate}

\subsection{AI Services and Predictive Modeling}
The system includes a dedicated AI service that instantaneously supplies market feedback to help farmers at harvest. The AI layer is an off-chain layer dedicated to predicting market prices, eliminating the need to pay high gas fees. Time-series algorithms are used to compare the farmers' attributes of input with market demand variables. The model predicts the latest price trends, allowing farmers to have a crystal-clear view of the monetary value of their crop before entering information into the ledger.

\section{AI/ML Price Prediction Pipeline}

\begin{figure}[htbp]
\centering
\includegraphics[width=0.95\linewidth]{farmchain_ai_flowchart.png}
\caption{The complete AI-Based Agricultural Price Prediction pipeline, detailing data collection, preprocessing, chronological splitting, XGBoost model training, and integration into the FastAPI microservice.}
\label{fig:ai_flowchart}
\end{figure}

This section describes the dataset, feature engineering, model selection, training methodology, and evaluation metrics for the agricultural price prediction component.

\subsection{Dataset}
Historical commodity price data was collected from three sources: the Government of India's Agmarknet API (data.gov.in), the CEDA Agmarknet portal, and the VegetableMarketPrice web service. The combined dataset contains \textbf{2{,}940 raw records} spanning 7 agricultural commodities (Brinjal, Cabbage, Carrot, Green Chilli, Onion, Potato, Tomato) across 7 major Indian wholesale markets (Bangalore, Chennai, Delhi, Hyderabad, Mumbai, Mysore, Pune) over a 60-day continuous window. Each record includes the date, commodity name, market name, minimum price, maximum price, modal price (target variable, in Rs/kg), and daily market arrivals (tonnes). A schema normalizer standardizes diverse source formats into a unified internal representation and converts Rs/Quintal to Rs/kg where the mean price exceeds 400. After per-group feature engineering and NaN removal (due to 30-day lag windows), \textbf{1{,}470 valid records} remain for model training.

\subsection{Feature Engineering}
We engineered \textbf{30 time-series features} without future data leakage using strict \texttt{shift(1)} operations on all rolling calculations. Features are grouped into six categories (Table~\ref{tab:features}).

\begin{table}[htbp]
\caption{Engineered Feature Groups for XGBoost Price Prediction}
\label{tab:features}
\centering
\resizebox{\linewidth}{!}{%
\begin{tabular}{lll}
\toprule
\textbf{Category} & \textbf{Features} & \textbf{Count} \\
\midrule
Price Lags & lag\_1, lag\_3, lag\_7, lag\_14, lag\_30 & 5 \\
Rolling Means & rolling\_mean\_3, \_7, \_14, \_30 & 4 \\
Rolling Std Dev & rolling\_std\_7, \_14, \_30 & 3 \\
Price Changes & abs \& pct changes at 1d, 7d, 30d & 6 \\
Calendar Signals & day\_of\_week, month, quarter, etc. & 6 \\
Arrival Dynamics & arrivals, arrival\_lag\_1, \_7, rolling\_mean\_7 & 4 \\
Categorical & commodity\_encoded, market\_encoded & 2 \\
\midrule
\textbf{Total} & & \textbf{30} \\
\bottomrule
\end{tabular}%
}
\end{table}

All lag and rolling features are computed using \texttt{shift(1)} to ensure that row $t$ contains only information available at time $t{-}1$, strictly preventing lookahead bias. Commodity and market are label-encoded to enable cross-market generalization.

\subsection{Model Selection and Training}
We selected XGBoost Regressor \cite{chen2016xgboost} due to its established effectiveness on tabular time-series data \cite{xiong2020price}. The model was configured with \texttt{n\_estimators}=100, \texttt{learning\_rate}=0.05, \texttt{max\_depth}=5, and \texttt{objective}=\texttt{reg:squarederror}.

The dataset was split \textbf{chronologically} (first 80\% for training, last 20\% for testing) to prevent temporal data leakage: training data spans Aug 13--Sep 5, 2026 (\textbf{1{,}176 rows}), and the test set spans Sep 6--Sep 11, 2026 (\textbf{294 rows}). No random shuffling was applied. Feature engineering was performed per commodity-market group independently before concatenation, preventing cross-group lag contamination.

\subsection{Baseline and Evaluation}
We implemented a \textbf{na\"ive persistence baseline} where the predicted price for day $t$ equals the observed price at day $t{-}1$ (i.e., $\hat{y}_t = y_{t-1}$). Table~\ref{tab:ml_metrics} compares both models across four standard regression metrics on the 294-sample chronological test set.

\begin{table}[htbp]
\caption{XGBoost vs.\ Na\"ive Baseline on Chronological Test Set (294 Samples)}
\label{tab:ml_metrics}
\centering
\begin{tabular}{lcccc}
\toprule
\textbf{Model} & \textbf{MAE} & \textbf{RMSE} & \textbf{MAPE} & \textbf{R\textsuperscript{2}} \\
 & (Rs/kg) & (Rs/kg) & (\%) & \\
\midrule
Na\"ive Persistence & 1.06 & 1.41 & 2.41 & 0.991 \\
\textbf{XGBoost} & \textbf{1.26} & \textbf{1.86} & \textbf{2.72} & \textbf{0.984} \\
\bottomrule
\end{tabular}
\end{table}

Both models achieve $R^2 > 0.98$, confirming that predicted prices explain over 98\% of actual market price variance across all commodity-market pairs. The na\"ive baseline shows marginally lower point-forecast error on next-day prediction, which is a well-documented phenomenon in commodity markets where short-term price persistence dominates \cite{xiong2020price}. However, the XGBoost model provides two critical advantages: (1) it captures non-linear relationships between arrivals volume, seasonal patterns, and price movements that enable more accurate \textit{multi-day} forecasting, and (2) it generalizes across all 49 commodity-market pairs via the encoded categorical features, eliminating the need for separate per-pair models.

For production deployment, the model performs recursive auto-regressive forecasting up to 7 days ahead. For each forecast step $i$: a synthetic row is created for date $t{+}i$, the feature engineer computes lag and rolling features using all available history including prior predictions, and XGBoost predicts the modal price. A retail markup multiplier (default $1.30\times$) estimates consumer-facing prices.

\section{Technical Implementation}

\subsection{Smart Contract Deployment}
The deployment of \texttt{FarmChain.sol} to the testnet of the Ethereum Sepolia network with the address \texttt{0xeE0e5471..a8b030472} was successful. We executed 9 completely successful transactions to register test harvests.

\begin{figure}[htbp]
\centering
\includegraphics[width=0.95\linewidth]{etherscan_sepolia.png}
\caption{Confirmed transactions on the public Sepolia Etherscan block explorer showing the live deployment of smart contracts and verified registerProduct transactions.}
\label{fig:etherscan}
\end{figure}

Product registration is initiated server-side through a backend relayer wallet, enabling a \textbf{wallet-less} farmer experience. The farmer never interacts with MetaMask or pays gas fees directly. The operational logic is given in Algorithm 1.

\begin{algorithm}
\caption{Crop Registration Logic (\texttt{registerProduct})}
\begin{algorithmic}[1]
\REQUIRE $msg.sender$ is an authenticated Farmer address
\STATE $timestamp \leftarrow \text{block.timestamp}$
\STATE $batchId \leftarrow \text{Keccak256}(msg.sender, timestamp, \text{randomNonce})$
\STATE $products[batchId] \leftarrow \text{Product}(msg.sender, timestamp)$
\STATE \text{emit } \text{HarvestRegistered}(batchId, msg.sender)
\RETURN $batchId$
\end{algorithmic}
\end{algorithm}

\subsection{FastAPI AI Microservice}
Our Python microservice exposes the primary production endpoint:
\begin{itemize}
    \item \textbf{\texttt{POST /api/farmer/price-prediction}}: Accepts commodity name, market, and forecast horizon (1--7 days). Queries Supabase for recent price history (with CSV fallback). Executes the recursive XGBoost forecasting loop. Returns current price, historical trend data, multi-day predictions, and a trend classification (\texttt{INCREASING}, \texttt{STABLE}, or \texttt{DECREASING}).
\end{itemize}
Successful predictions are published as events to the Kafka \texttt{price-predictions} topic for downstream consumption.

\subsection{Backend API \& Database}
The Node.js and Express server implements role-gated API routes:
\begin{itemize}
    \item \textbf{Auth Routes (\texttt{/api/auth}):} Registers, hashes passwords securely using bcrypt and generates JWTs. Role-based approval queues ensure pending roles await admin confirmation.
    \item \textbf{Product Routes (\texttt{/api/products}):} Stores complex data (such as crop names, high-resolution image URLs) directly into our Supabase database. On product creation, the system calls the AI microservice for quality scoring and queues a \texttt{BATCH\_CREATED} event to the Kafka outbox.
    \item \textbf{Admin Routes (\texttt{/api/admin}):} User approval, product approval (triggering on-chain \texttt{registerProduct}), system analytics, and blockchain health monitoring.
    \item \textbf{Retailer Routes (\texttt{/api/retailer}):} Purchase requests, shipment tracking, QR-based batch verification, on-chain receipt recording, inventory management, POS sales, wastage reporting, and customer returns.
\end{itemize}

\subsection{Frontend Dashboards}
Our farmer interface was built with React 18, Tailwind CSS, Lucide for clear and responsive icons, and the Ethereum blockchain interface Ethers.js to connect the old-fashioned web interface with the blockchain.

After successful login, users will be shown the \textbf{Farmer Dashboard}. At the top of the screen are three large summary cards: Total Batches, Harvested Batches, Volume Shipped. Under these metrics, a high resolution picture of the crop and organic tags and logged expiry dates are shown on the visual inventory cards. If the ``Add Harvest'' button is clicked, a simple pop-up form appears for new registrations.

Beyond the farmer, the platform provides dedicated dashboards for each downstream role:
\begin{itemize}
    \item \textbf{Processor Dashboard:} A comprehensive component managing gate-scale verification, multi-criteria quality inspection (Appearance, Freshness, Cleanliness scoring with Accept/Reject decisions), batch processing stages, and cold-chain distributor handoff.
    \item \textbf{Distributor Dashboard:} Blockchain/batch ID lookup, custody intake via \texttt{transferOwnership}, transit management, and delivery performance analytics.
    \item \textbf{Retailer Suite:} Eleven dedicated pages covering QR-code product receipt, inventory management, point-of-sale checkout, wastage tracking, expiry monitoring, customer returns, and full batch traceability.
    \item \textbf{Consumer Tracker:} Public \texttt{/track} page with triple-layer QR scanning (live camera, \texttt{BarcodeDetector}, \texttt{jsQR}), direct on-chain verification via \texttt{ethers.JsonRpcProvider}, and a six-milestone authenticity checklist (Registration, Farmer Origin, Processor Inspection, Distributor Logistics, Retailer Delivery, Blockchain Record).
\end{itemize}

\section{Experimental Evaluation}

\subsection{Smart Contract Execution \& Gas Metrics}
We stress-tested the gas costs of this decentralized approach directly on the Sepolia network (Table \ref{tab:gas_metrics}) to assess its practicality for an average farmer.

\begin{table}[htbp]
\caption{FarmChain Smart Contract Execution \& Gas Metrics on Sepolia}
\label{tab:gas_metrics}
\centering
\resizebox{\linewidth}{!}{%
\begin{tabular}{llrr}
\toprule
\textbf{Function} & \textbf{State Change} & \textbf{Gas Used} & \textbf{Fee (ETH)} \\
\midrule
\texttt{Deployment} & Bytecode deployment & $\approx 850{,}000$ & $\approx 0.00170$ \\
\texttt{registerProduct()} & Product struct + event & $\approx 185{,}432$ & $\approx 0.00037$ \\
\texttt{updateStage()} & Status + history append & $\approx 65{,}218$ & $\approx 0.00013$ \\
\texttt{transferOwnership()} & Owner reassign + log & $\approx 70{,}104$ & $\approx 0.00014$ \\
\texttt{getProductHistory()} & Read-only query & 0 & Free \\
\bottomrule
\end{tabular}%
}
\end{table}

According to our in-depth analyses of the transactions made on Etherscan, it takes around 185,432 units of computational gas to register a brand new crop batch onto the ledger. The cumulative gas cost for a complete farm-to-fork journey (register + 3 stage updates + 3 ownership transfers) totals approximately 0.00120 ETH, which is negligible compared to enterprise licensing fees. Moreover, free read calls are inherent to blockchain architecture, which means that the farmer can use blockchain for free to verify their historical product logs.

\subsection{System Latency}
The Node.js API typically takes under an astounding 14ms to respond. The average time for the Python AI endpoints is about 35 milliseconds. The total time for all these networked operations - from the exact moment the farmer clicks on the last submit button to the completion of the AI price predictions creation, writing data to the database, and updating the localized dashboard - is less than 1.5 seconds.

\subsection{Workflow Testing}
To test the stability of the whole system, we ran several dozen concurrent transactions on the frontend UI across all six roles. The mathematical answers of the AI microservice have been successfully verified and the final transaction hashes for the cryptographic transactions have been manually cross-referenced on the Sepolia Etherscan explorer to guarantee 100\% data fidelity. The smart contract's \texttt{require(currentOwner == msg.sender)} guard successfully rejected all unauthorized stage-update and ownership-transfer attempts during testing.

\section{Discussion}

\subsection{Key Findings}
\begin{enumerate}
    \item \textbf{Complete Traceability:} The prototype demonstrates end-to-end farm-to-fork traceability with verifiable custody transfers at every stage, not merely farmer-side harvest logging.
    \item \textbf{Actionable AI:} An XGBoost model achieving $R^2 = 0.984$ and MAPE of 2.72\% provides farmers with genuinely useful market intelligence, reducing information asymmetry against middlemen.
    \item \textbf{Off-Chain Validation:} The three-gate validation architecture (admin approval, processor physical inspection, AI anomaly detection) addresses the fundamental limitation that blockchain immutability cannot guarantee input truthfulness.
    \item \textbf{Decoupled Performance:} Using the intensive AI mathematical models as a standalone Python microservice \cite{dragoni2017microservices} kept the main React website extremely responsive \cite{zheng2018blockchain}.
\end{enumerate}

\subsection{Limitations \& Next Steps}
The current XGBoost model was trained on 2{,}940 records spanning 60 days. Expanding the dataset to multi-year horizons would improve seasonal pattern capture. The AI endpoints for quality classification, disease detection, and shelf-life prediction currently use heuristic-based scoring and represent targets for future deep-learning replacement. We will also scale the underlying smart contracts on a Layer-2 scaling network, like the Polygon Amoy network, in the future. This will keep the gas fees the mainnet transaction users pay close to zero, so that the platform is always affordable for farmers with limited financial resources \cite{rejeb2019leveraging}.

\section{Conclusion}
The full development process of an open-source, agricultural tracking platform, called \textbf{FarmChain AI} was described in this paper. The platform implements the complete supply chain lifecycle across six stakeholder roles with verifiable on-chain custody transfers. The XGBoost price prediction model, trained on 2{,}940 records with 30 leakage-free features and evaluated via chronological splitting, achieves a MAPE of 2.72\% and $R^2$ of 0.984. Multi-gate off-chain validation addresses input truthfulness concerns. After thorough testing, we found that the integrated system has an extremely high speed and is very low cost when running on the blockchain ($\approx$0.00037 ETH per registration). In conclusion, the project's results show that a secure, intelligent and transparent agricultural logging system can be implemented and function effectively in the interests of empowering farmers to access the benefits of logging without the need to rely on the big budgets of large-scale corporate entities.

\section*{Acknowledgment}
The authors would like to thank the faculty and students of the Dept. of Computer Science \& Engineering and Dept. of Information Science \& Engineering, RNS Institute of Technology, Bengaluru, India for their guidance, support and encouragement.

\begin{thebibliography}{00}

\bibitem{india_agri_stats} Ministry of Agriculture \& Farmers Welfare, Government of India, ``Agricultural Statistics at a Glance 2023,'' 2023.

\bibitem{farmer_income_study} National Bank for Agriculture and Rural Development (NABARD), ``Report on Financial Inclusion and Agricultural Development in India,'' 2022.

\bibitem{tian2017agrifood} F. Tian, ``An agri-food supply chain traceability system for China based on RFID \& blockchain technology,'' in \textit{Proc. 13th Int. Conf. Service Systems and Service Management (ICSSSM)}, 2016, pp. 1--6.

\bibitem{lin2019bifts} Q. Lin et al., ``Food Safety Traceability System Based on Blockchain and EPCIS,'' \textit{IEEE Access}, vol. 7, pp. 20698--20707, 2019.

\bibitem{kamath2018food} R. Kamath, ``Food Traceability on Blockchain: Walmart's Pork and Mango Pilots with IBM,'' \textit{The Journal of the British Blockchain Association}, vol. 1, no. 1, 2018.

\bibitem{salah2019soybean} K. Salah, N. Nizamuddin, R. Jayaraman, and M. Omar, ``Blockchain-based soybean traceability in agricultural supply chain,'' \textit{IEEE Access}, vol. 7, pp. 73295--73305, 2019.

\bibitem{antonucci2019review} F. Antonucci et al., ``A review on blockchain applications in the agri-food sector,'' \textit{Journal of the Science of Food and Agriculture}, vol. 99, no. 14, pp. 6129--6138, 2019.

\bibitem{kamilaris2018deep} A. Kamilaris and F. X. Prenafeta-Bold\'u, ``Deep learning in agriculture: A survey,'' \textit{Computers and Electronics in Agriculture}, vol. 147, pp. 70--90, 2018.

\bibitem{liakos2018ml} K. G. Liakos et al., ``Machine Learning in Agriculture: A Review,'' \textit{Sensors}, vol. 18, no. 8, 2018.

\bibitem{xiong2020price} Y. Xiong et al., ``Agricultural product price prediction based on machine learning,'' \textit{IEEE Access}, vol. 8, pp. 71505--71513, 2020.

\bibitem{chen2016xgboost} T. Chen and C. Guestrin, ``XGBoost: A Scalable Tree Boosting System,'' in \textit{Proc. 22nd ACM SIGKDD}, 2016, pp. 785--794.

\bibitem{wood2014ethereum} G. Wood, ``Ethereum: A secure decentralised generalised transaction ledger,'' \textit{Ethereum Project Yellow Paper}, vol. 151, pp. 1--32, 2014.

\bibitem{zheng2018blockchain} Z. Zheng, S. Xie, H. N. Dai, X. Chen, and H. Wang, ``Blockchain challenges and opportunities: A survey,'' \textit{Int. Journal of Web and Grid Services}, vol. 14, no. 4, pp. 352--375, 2018.

\bibitem{dragoni2017microservices} N. Dragoni et al., ``Microservices: yesterday, today, and tomorrow,'' in \textit{Present and Ulterior Software Engineering}, Springer, 2017, pp. 195--216.

\bibitem{atzei2017survey} N. Atzei, M. Bartoletti, and T. Cimoli, ``A survey of attacks on Ethereum smart contracts (SoK),'' in \textit{Int. Conf. on Principles of Security and Trust}, 2017, pp. 164--186.

\bibitem{rejeb2019leveraging} A. Rejeb, J. G. Keogh, and H. Treiblmaier, ``Leveraging the Internet of Things and Blockchain Technology in Supply Chain Management,'' \textit{Future Internet}, vol. 11, no. 7, p. 161, 2019.

\end{thebibliography}

\end{document}
